\scanpage{107}

\hfill 21.2.\\

\handwritten{8.} Returning to theamalgamated sum $X=\handwritten{\bigcup_i}X_i$, the cellular decompositions of the components $X_i\handwritten{\simeq}D_{n(i)}$ define a cellular decomposition of $X$, whose set of cells with incidence relation forms a finite ordered set $K$, finite union of a family of subsets $(K_i)_{i\handwritten{\in}I}$, with given isomorphisms\deleted{\unsure{ of}}
\[
f_i:K_i\handwritten{\simeq}\underline{J}_{n(i)}\qquad(i\handwritten{\in}I)\quad ,
\]
where for every index $n\handwritten{\in}\underline{N}$, $\underline{J}_n$ denotes the ordered set of the $2n+1$ cells $S_j^{\handwritten{\varepsilon}}$ $(0\handwritten{\leq}j\handwritten{\leq}n-1,\handwritten{\varepsilon\in\{}\pm1\handwritten{\}}),D_n$ of \deleted{\unsure{$D_n$}}the pertinent cellular decomposition of $D_n$. We may without loss of generality assume there is no\deleted{\unsure{n}} inclusion relation between the $K_i$, moreover the standartness condition described above readily translates into a condition on this structure $K$, $(K_i)_{i\handwritten{\in}I}$, $(f_i)_{i\handwritten{\in}I}$, and implies that for $i,i'\handwritten{\in}I$, $K_i\handwritten{\to}K_{i'}$ is a``closed'' subset in the two ordered sets $K_i,K_{i'}$ (namely contains with any element $x$ the elements smaller than $x$), and moreover isomorphic (for the induced order) to some $\underline{J}_n$. Thus the category $B_{\unsure{o|0}}$ \deleted{is} can be viewed as the category of such ``standart ordered sets'' (with the extra structure on these just said), and the category $C_o$ can be defined most simply as the dual category $B_o^{o}$. (NB the definition of morphisms in $B_o$ is clear I guess \ldots) I believe the category $B_o$ is stable under amalgamated sums $X\handwritten{\amalg}_ZY$, provided \inserted{however} we insist that the empty structure $X$ is \uline{not} allowed -- otherwise we have to \deleted{insist} restrict to amalgamated sums with $Z\handwritten{\ne\varnothing}$. It seems finally more convenient to exclude the empty structure in $B_o$, i.e. to exclude final element from $\underline{C}_o$, for the benefit of being able to state that $\underline{C}_o$ (and \uline{all categories} $\underline{C}_n$\handwritten{*}) \deleted{should be} are stable under binary amalgamated sums, and that the functors $\underline{C}_n\handwritten{\to}C$ (and ultimately $\underline{C}_{\handwritten{\infty}}\handwritten{\to}C$) we are interested in are just those commuting to arbitrary binary amalgamated sums (without \deleted{\unsure{refere}} awkwa\handwritten{\unsure{rd}} reference to the objects $\underline{F}_i$ and the iterated source and target maps between them \ldots).
% slanted
\margin{\handwritten{This is obviously false, cf \unsure{PS|P.S} p.12}}

\handwritten{9.} The category $\underline{C}_o$ being fairly well understood, remains to complete the construction by the inductive step, passing from $\underline{C}_n$ to $\underline{C}_{n+1}$. The main properties I have in mind therefore, for the sequence of categories $\underline{C}_n$ and their limit $\underline{C}_{\handwritten{\infty}}$, are the following two.

A) For any $\unsure{K}\handwritten{\in}\operatorname{Ob}(\underline{C}_{\handwritten{\infty}})$ $(\unsure{=}\operatorname{Ob}(\underline{C}_o))$, and any two arrows in $\underline{C}_{\handwritten{\infty}}$
\[
 f,g: \quad K\handwritten{\rightrightarrows}F_i\quad ,
\]
with $i\handwritten{\in}\underline{N}$\deleted{\unsure{,}}, and such that either $i=0$, or the equalities
\[
 s_{\unsure{1|l}}^i f=s_1^i g \quad , \quad t_1^i f=t_1^i g \tag*{\handwritten{(1)}}
\]
hold (case $i\handwritten{\geq}1$), there exists $h:K\quad F_{i+1}$ such that
\[
 s_1^{i+1}h=f\quad ,\quad t_1^{i+1}h=g\quad . \tag*{\handwritten{(2)}}
\]

B) For any $n\handwritten{\in}\underline{N}$, the category $\underline{C}_{n+1}$ is deduced from $\underline{C}_n$ by keeping the same objects, and just adding new arrows \deleted{\unsure{as}} $h$ as in A), with $f,g$ arrows in $\unsure{C|\,\underline{C}_n}$.

% archive sheet number: 106
